\documentclass[9pt,twocolumn]{extarticle}

\usepackage[top=0.60in,bottom=0.62in,left=0.58in,right=0.58in,columnsep=0.24in]{geometry}
\usepackage[T1]{fontenc}
\usepackage{lmodern}
\usepackage{microtype}
\usepackage[table]{xcolor}
\usepackage{booktabs}
\usepackage{tabularx}
\usepackage{enumitem}
\usepackage{amsmath}
\usepackage{amssymb}
\usepackage{fancyhdr}
\usepackage[compact]{titlesec}
\usepackage[hidelinks]{hyperref}

\definecolor{DDRAMBlue}{HTML}{4773B4}
\definecolor{DDRAMGreen}{HTML}{59AA9B}
\definecolor{DDRAMOrange}{HTML}{FF9D42}
\definecolor{DDRAMViolet}{HTML}{6451DE}
\definecolor{DDRAMGrey}{HTML}{ECECEC}
\definecolor{DDRAMCharcoal}{HTML}{4B514F}
\definecolor{DDRAMDeferred}{HTML}{777777}

\hypersetup{
  colorlinks=true,
  linkcolor=DDRAMBlue,
  urlcolor=DDRAMViolet,
  pdftitle={DDRAM-O Milestone 1 Specification},
  pdfauthor={DDRAM Project}
}

\pagestyle{fancy}
\fancyhf{}
\lhead{DDRAM-O Milestone 1}
\rhead{Simplified Device and Timing Model}
\cfoot{\thepage}
\setlength{\headheight}{14pt}

\setlength{\parindent}{0.9em}
\setlength{\parskip}{0pt}
\setlength{\textfloatsep}{7pt plus 2pt minus 2pt}
\setlength{\floatsep}{6pt plus 2pt minus 2pt}
\setlength{\intextsep}{6pt plus 2pt minus 2pt}
\setlength{\abovecaptionskip}{4pt}
\setlength{\belowcaptionskip}{2pt}
\renewcommand{\arraystretch}{0.94}
\setlist{nosep,leftmargin=*,topsep=2pt}
\titlespacing*{\section}{0pt}{7pt plus 2pt minus 1pt}{3pt}
\titlespacing*{\subsection}{0pt}{5pt plus 1pt minus 1pt}{2pt}
\titleformat{\section}{\large\bfseries\color{DDRAMCharcoal}}{\thesection}{0.45em}{}
\titleformat{\subsection}{\normalsize\bfseries\color{DDRAMCharcoal}}{\thesubsection}{0.4em}{}
\newcommand{\cmd}[1]{\texttt{#1}}
\newcommand{\param}[1]{\texttt{#1}}
\newcommand{\viewA}{\textsf{A}}
\newcommand{\viewB}{\textsf{B}}
\newcommand{\provisional}{\textcolor{DDRAMOrange}{\textbf{Provisional}}}
\newcommand{\normative}{\textcolor{DDRAMBlue}{\textbf{Normative}}}
\newcommand{\timingrevision}{M1-P1}
\newcommand{\deferredcell}[1]{\textcolor{DDRAMDeferred}{#1}}

\begin{document}
\twocolumn[
\begin{center}
  {\LARGE\bfseries DDRAM-O Milestone 1 Specification}\\
  \vspace{0.2em}
  {\large Simplified Orthogonal-View Device and Timing Model}\\
  \vspace{0.35em}
  {\small DDRAM Project \hfill Revision 0.7 --- 21 October 2026}
\end{center}
\vspace{2pt}
\noindent\fcolorbox{DDRAMBlue}{DDRAMGrey}{%
\parbox{\dimexpr\textwidth-2\fboxsep-2\fboxrule\relax}{%
\textbf{Status.} This document defines the first implementable DDRAM-O model
for Ramulator2. Architectural interfaces and state transitions are normative.
Electrical timing values introduced specifically for DDRAM-O are placeholders
and must remain configurable until replaced by SPICE-calibrated values.}}
\vspace{7pt}
]

\tableofcontents

\section{Purpose and Scope}

DDRAM-O provides two access structures over the same DRAM cell array. View
\viewA{} uses the conventional row orientation; view \viewB{} is rotated by
$90^\circ$, so a row in one view geometrically corresponds to a column in the
other. The complete architecture creates cell-intersection hazards between
simultaneously open rows. Milestone 1 intentionally omits those hazards and
models the two row paths as independent, subject to an explicit cross-view
activation delay.

The purpose of this milestone is to establish a stable device interface for:

\begin{itemize}
  \item view-specific conventional accesses;
  \item overlap between activity on the two row paths;
  \item atomic processing-using-DRAM (PuD) primitives;
  \item configurable, visible placeholder timings; and
  \item command traces suitable for correctness and scheduling analysis.
\end{itemize}

This is a timing and state model. Functional arithmetic, operand-row
allocation, energy calibration, and operation scheduling are later milestones.

\subsection{Normative language}

The words \textbf{must}, \textbf{must not}, \textbf{should}, and \textbf{may}
describe implementation requirements. A requirement marked \provisional{}
defines the initial model but may be replaced by calibrated data without
changing the command interface.

\section{Goals and Non-Goals}

\subsection{Milestone 1 goals}

\begin{enumerate}
  \item Introduce a distinct \cmd{DDRAM\_O} DRAM standard.
  \item Preserve the conventional DDR5 access path through view \viewA{}.
  \item Permit one conventional row to remain open independently in each view.
  \item Introduce explicit activation, access, and precharge commands per view.
  \item Introduce an atomic multi-row activation primitive per view.
  \item Permit cross-view overlap after a configurable minimum delay.
  \item Preserve conventional command-, column-, and data-bus sharing.
  \item Make every DDRAM-O-specific timing overridable from Python configuration.
\end{enumerate}

\subsection{Explicit non-goals}

Milestone 1 does \emph{not} model:

\begin{itemize}
  \item A/B row-intersection conflicts or simultaneous-sensing hazards;
  \item masked intersection bits or half-amplitude sensing;
  \item persistent states with three or five conventionally open rows;
  \item functional data values or arithmetic correctness;
  \item compute-row allocation, operand lifetime, or spilling;
  \item calibrated multi-row activation current or energy; or
  \item end-to-end CPU trace offloading.
\end{itemize}

Results produced by this model must not be described as SPICE-calibrated.

\section{Architectural Model}

\subsection{Views and addressing}

View \viewA{} is the host-oriented conventional view. View \viewB{} is the
orthogonal view. Both commands use the existing hierarchical Ramulator address
vector. In Milestone 1, the row field is interpreted in the coordinate system
of the selected view. Consequently, row number $r$ in \viewA{} and row number
$r$ in \viewB{} do not denote the same geometric row.

More precisely, the row field is a \emph{view-local line index}. For example,
\cmd{A:10} names line 10 in the \viewA{} coordinate system, whereas
\cmd{B:10} names line 10 in the orthogonal \viewB{} coordinate system. The
equal integers are labels in different namespaces; they do not mean that one
physical row has been opened twice. The two selected lines may physically
intersect at a cell, but Milestone 1 assigns that intersection no additional
state transition, timing penalty, masking rule, or data-value effect. Thus
both activations are legal in the simplified timing model. A later
intersection-aware functional model must determine the shared-cell mapping and
any resulting electrical or correctness restrictions.

The address mapper is not required to perform a physical transpose in this
milestone. The selected command identifies the coordinate system. A later
functional model will map both coordinate systems onto shared cells.

\subsection{Data layout and operation orientation}
\label{sec:layout}

The initial controller integration shall use view \viewA{} for conventional
host accesses and view \viewB{} for bit-serial PuD. These are the baseline
layout contracts, not interchangeable row-buffer placement choices.
Bit-serial computation processes successive bit positions over time while
operating on many independent values in parallel across sensing lanes.
An operand's bit planes must therefore occupy the view-local lines expected
by the PuD primitive, with matching lane identities for all operands.

For illustration, let $C[i,j]$ denote cells in a rectangular tile, where an
A-line fixes $i$ and a B-line fixes $j$. A compatible placement may store
the bits of value $x_i$ along A-line $i$:
\begin{equation*}
 C[i,b]=x_i[b],\qquad B_b=(x_0[b],x_1[b],\ldots).
\end{equation*}
Here $B_b$ exposes a bit plane: lane $i$ retains the same value identity as
the bit index $b$ advances. With several operands, separate B-lines hold
their respective bit planes; a primitive combines corresponding lanes.
This is an illustrative placement, not a fixed DDR5 address decoder or an
assumption that one host value occupies an entire physical A-row.

Orthogonal access can expose an already compatible placement without an
explicit transpose. It does not guarantee that arbitrary host buffers are
compatible: packing, bit order, operand placement, tile geometry, and lane
alignment must agree with the operation. Changing \cmd{selected\_view} does
not transpose, copy, or repack data. A layout mismatch requires an explicitly
modeled conversion or rejection before execution; conversion costs must not
be hidden in view selection.

Low-level \cmd{MRA\_A} remains a supported device command. It does not make
an arbitrary bit-serial PuD operation valid in A. Other operation/layout
combinations may permit A only when their mapping and primitive semantics
are explicitly defined. Likewise, \cmd{RD\_B}/\cmd{WR\_B} address B-local
lines; they are not substitutes for host accesses to the same numeric A-row.

\subsection{Per-bank state}

Each physical bank maintains the state shown in Table~\ref{tab:state}.

\begin{table}[ht]
\centering
\caption{Required simplified state per physical bank.}
\label{tab:state}
\rowcolors{2}{DDRAMGrey}{white}
\begin{tabularx}{\columnwidth}{@{}l l X@{}}
\toprule
Field & Domain & Meaning \\
\midrule
\cmd{active\_row[A]} & closed or row ID & Conventional row open through view \viewA{}. \\
\cmd{active\_row[B]} & closed or row ID & Conventional row open through view \viewB{}. \\
\cmd{busy\_until[A]} & cycle & Earliest cycle at which another \viewA{} row-path operation may start. \\
\cmd{busy\_until[B]} & cycle & Earliest cycle at which another \viewB{} row-path operation may start. \\
\bottomrule
\end{tabularx}
\end{table}

The two active-row fields are independent. Opening or precharging one view must
not change the other view's state. No equality check between the two row IDs is
performed. A bank is conventionally closed only when both fields are closed.

\subsection{Simplified concurrency rule}

Commands on different views may overlap once their activation-like command
starts are separated by \param{nXACTO}. Commands still contend for all shared
resources inherited from DDR5. In particular, DDRAM-O adds row-path
parallelism, not an additional off-chip data interface.

\section{Command Interface}

\subsection{Required commands}

Table~\ref{tab:commands} defines the Milestone 1 command set. ``Local'' means
the state or timing effect applies only to the command's selected view.

\begin{table*}[t]
\centering
\caption{Required DDRAM-O commands.}
\label{tab:commands}
\begin{tabularx}{\textwidth}{@{}p{0.14\textwidth}p{0.13\textwidth}X@{}}
\toprule
Command & Scope & Semantics \\
\midrule
\rowcolor{DDRAMGrey}\cmd{ACT\_A} & Bank, local & Open the target row in view \viewA{}. \\
\cmd{ACT\_B} & Bank, local & Open the target row in view \viewB{}. \\
\rowcolor{DDRAMGrey}\cmd{PRE\_A} & Bank, local & Close only the row open in view \viewA{}. \\
\cmd{PRE\_B} & Bank, local & Close only the row open in view \viewB{}. \\
\rowcolor{DDRAMGrey}\cmd{RD\_A}, \cmd{WR\_A} & Bank, shared data path & Read or write the row open in view \viewA{}. \\
\cmd{RD\_B}, \cmd{WR\_B} & Bank, shared data path & Read or write the row open in view \viewB{}. \\
\rowcolor{DDRAMGrey}\cmd{MRA\_A} & Bank/subarray, local & Execute an atomic multi-row activation on view \viewA{}; no persistent multi-row-open state remains. \\
\cmd{MRA\_B} & Bank/subarray, local & Symmetric atomic multi-row activation on view \viewB{}. \\
\rowcolor{DDRAMGrey}\cmd{PREab} & Rank & Close both views in all targeted banks, for maintenance compatibility. \\
\cmd{PREsb} & Same-bank scope & Close both views in the targeted banks, for maintenance compatibility. \\
\bottomrule
\end{tabularx}
\end{table*}

\subsection{Multi-row activation abstraction}

\cmd{MRA\_A} and \cmd{MRA\_B} represent the complete sequence

\begin{equation}
  \text{multi-row activation}
  \rightarrow \text{sense/majority}
  \rightarrow \text{restore}
  \rightarrow \text{precharge}.
\end{equation}

The command must begin with its local view closed. If that view is open, its
prerequisite is the corresponding local precharge. The opposite view may
remain open. At command issue, the local view enters a timed busy state but
does not acquire a persistent active-row identity. It is logically closed
after the atomic operation.

Milestone 1 may use the command address as a representative bank/subarray
target. The operand-row list belongs to the compound PuD request introduced in
Milestone 2 and must not be encoded through unrelated address-vector fields.

\subsection{Reserved primitive timings}

The implementation should also define \param{nCOPY} and \param{nNOT} timing
parameters now, even if explicit \cmd{COPY\_A/B} and \cmd{NOT\_A/B} commands
are introduced with functional primitives in a later milestone. Reserving
separate parameters prevents future calibration from coupling logically
different primitives merely because their initial values are equal.

\subsection{PuD command and operation reference}
\label{sec:pud-catalogue}

Table~\ref{tab:pud-catalogue} is the reference for planned PuD commands and
their layout-to-view mapping. Device primitives and compiler-level operations
are distinguished explicitly: an operation name does not imply a new DRAM
opcode. Layout compatibility follows Section~\ref{sec:layout}; timing symbols
refer to Table~\ref{tab:parameters} and the profiles in
Section~\ref{sec:pumulator}. The timing ledger remains \timingrevision{}.

\textbf{Initial implementation target: PuD executes through view B;
ordinary host accesses use view A.} The highlighted B primitives define the
initial PuD controller scope. Greyed A-view PuD primitives are deferred,
optional extensions and must not be generated by the initial PuD compiler
or selected by the controller. Existing low-level \cmd{MRA\_A} device
support and device tests remain valid. This deferral does not restrict
ordinary A-view reads, writes, or their prerequisites.

\begin{table*}[t]
\centering
\small
\caption{PuD command/operation catalogue. Blue B-primitive rows are the initial
PuD target; grey A-primitive rows are deferred optional extensions.
Timings are provisional, not calibrated.}
\label{tab:pud-catalogue}
\rowcolors{2}{DDRAMGrey}{white}
\begin{tabularx}{\textwidth}{@{}p{0.15\textwidth}p{0.25\textwidth}p{0.05\textwidth}p{0.16\textwidth}X@{}}
\toprule
Command / operation & Layout and operand requirement & View & Service time & Status / remaining requirements \\
\midrule
\rowcolor{DDRAMBlue!10}\cmd{MRA\_B} & Bit-serial bit planes on B-local operand lines; matching value lanes. & B & \param{nMRA} & Timing-only device and controller path implemented, including service completion. Operand lists, physical layout validation, and functional majority pending. \\
\rowcolor{DDRAMBlue!10}\cmd{COPY\_B} & Source and destination are B-local bit planes with identical lane and bit interpretation in the same local compute domain. & B & \param{nCOPY} & Initial PuD target; timing parameter reserved. Requires operand encoding, destination placement, and copy-path support. No transpose implied. \\
\rowcolor{DDRAMBlue!10}\cmd{NOT\_B} & B-local source and destination with matching lanes; a DCC-compatible inversion destination/path. & B & \param{nNOT} & Initial PuD target; timing parameter reserved. Copy-equivalent default assumes DCC support; functional path pending. \\
\rowcolor{DDRAMGrey}\deferredcell{\cmd{MRA\_A}} & \deferredcell{Explicit A-local operand-line mapping for a supported operation; not the baseline bit-serial layout.} & \deferredcell{A} & \deferredcell{\param{nMRA}} & \deferredcell{Deferred for PuD integration. Low-level device timing/state exists; functional A-view PuD is an optional extension.} \\
\rowcolor{DDRAMGrey}\deferredcell{\cmd{COPY\_A}} & \deferredcell{A-local source/destination with matching packed-bit interpretation in the same local compute domain.} & \deferredcell{A} & \deferredcell{\param{nCOPY}} & \deferredcell{Deferred optional extension; no opcode implemented. Requires a defined A-local copy path and operands.} \\
\rowcolor{DDRAMGrey}\deferredcell{\cmd{NOT\_A}} & \deferredcell{A-local source/destination with matching bit interpretation and a supported DCC inversion path.} & \deferredcell{A} & \deferredcell{\param{nNOT}} & \deferredcell{Deferred optional extension; no opcode implemented. Requires independently confirmed A-local inversion support.} \\
MAJ3 bit-plane operation & Three aligned B-local operand bit planes, temporary compute lines, and one destination bit plane. & B & $3\param{nCOPY}$\newline${}+\param{nMRA}$\newline${}+\param{nCOPY}$ & Planned compiler sequence, not an opcode. Baseline preserves three inputs and copies one result; elision requires proven placement/liveness. \\
Addition / multiplication & Compatible bit-serial operands and intermediates, declared width/bit order, aligned lanes, and allocated temporary lines. & B & Dependency-graph cost; no fixed scalar timing & Planned compound operations. Decompose into MRA/COPY/NOT leaves; include carry/intermediate dependencies and required spills. \\
Layout conversion / cross-view movement & Explicit source and destination layout descriptors, coordinate mapping, and dependency ordering. & A/B & Unspecified; no timing parameter yet & Future explicit mechanism, not an implemented opcode. May be unnecessary for an already compatible placement; otherwise conversion is not free. \\
\bottomrule
\end{tabularx}
\end{table*}

The initial model uses the same symbolic primitive cost in A and B; this is
a provisional symmetry assumption, not evidence of equal circuit latency.
At DDR5\_3200AN under the default split profile, the reference values are
\param{nMRA}$=76$, \param{nCOPY}$=83$, and \param{nNOT}$=83$ cycles.
The MAJ3 bit-plane baseline is consequently 408 cycles; it is not a complete
multi-bit addition or multiplication. Explicit overrides and alternate copy
profiles supersede these example numbers.

Service time is measured from primitive issue and includes the internal final
precharge assumed by AP/AAP. A prerequisite precharge of an initially open
view is additional. Shared bus/power limits, cross-view \param{nXACTO},
dependencies, and contention can delay issue or completion further. MRA's
command/address occupancy remains two cycles, distinct from \param{nMRA};
COPY/NOT command encoding, occupancy, and prerequisite rules still require
definition. They must join applicable activation constraints when implemented.
Ordinary host and B-local read/write commands retain the DDR5 timing rules in
Tables~\ref{tab:commands} and~\ref{tab:constraints}.

For a serialized graph with no additional movement, the primitive service sum
is
\begin{equation*}
 L_{\mathrm{seq}}=N_M\param{nMRA}+N_C\param{nCOPY}+N_N\param{nNOT},
\end{equation*}
where the counts refer to emitted leaves after placement and spill insertion.
Count each inversion either as one NOT leaf or as its expanded copy sequence,
never both. This sum is a service-cost reference; actual scheduled latency
depends on graph dependencies and resource availability.

When a command or layout requirement changes, update this catalogue, its
mapping/validation rules, and the corresponding acceptance tests together.
Keep formulas in the timing ledger as the source of defaults; regenerate
numeric examples for the selected preset after any timing/profile change.
Increment the timing-ledger revision when timing assumptions change, and the
document revision for catalogue or interface-only changes. Every future
primitive must specify layout/operands, legal view, prerequisites, service and
CA occupancy, completion semantics, and implementation status before it is
enabled by the controller.

\section{Timing Model}

\subsection{Parameters and initial values}

Command delays are positive integer DRAM cycles; the decoder-delay input
$t_{\mathrm{DEC}}$ is specified in nanoseconds. Table~\ref{tab:parameters}
distinguishes Pumulator-inherited primitive assumptions from DDRAM-O-specific
placeholders. The current
provisional timing ledger is revision \timingrevision{}. This identifier and
the resolved values should accompany experiment results until calibrated
parameters replace the ledger.

\begin{table*}[t]
\centering
\caption{DDRAM-O-specific timing parameters.}
\label{tab:parameters}
\rowcolors{2}{DDRAMGrey}{white}
\begin{tabularx}{\textwidth}{@{}l l l X@{}}
\toprule
Parameter & Initial value & Status & Interpretation \\
\midrule
\param{nXACTO} & $8$ & \provisional{} & Minimum start-to-start delay between activation-like commands on opposite views of one bank/subarray. \\
\param{nMRA} & $\param{nRAS}+\param{nRP}$ & \provisional{} & Local-view occupancy of an atomic multi-row activation, including restoration and precharge. \\
\param{nCOPY} & $\param{nRAS}+\param{nRP}+n_{\mathrm{DEC}}$ & \provisional{} & Pumulator split-decoder activate--activate--precharge (AAP) copy cost; $n_{\mathrm{DEC}}=\lceil t_{\mathrm{DEC}}/t_{\mathrm{CK}}\rceil$, initially $t_{\mathrm{DEC}}=4$ ns. \\
\param{nNOT} & $\param{nCOPY}$ & \provisional{} & Initial inversion cost when represented as a DCC-assisted copy. \\
\bottomrule
\end{tabularx}
\end{table*}

The default \param{nXACTO} equals the current provisional DDRAM-P cross-view
delay for comparability; it is not an electrical claim. The \param{nMRA} and
\param{nCOPY} defaults inherit Pumulator's SIMDRAM AP/AAP cost model, with
unit base and primitive scaling factors and ceiling conversion to cycles.
This inheritance supplies an interim modeling baseline, not electrical
validation for DDRAM-O. In particular, the copy default assumes a split row
decoder and \param{nNOT} assumes a dual-contact-cell (DCC) inversion path;
these circuit capabilities must be confirmed for the eventual device.
SPICE calibration may replace the values without renaming commands.

\subsection{Pumulator provenance and copy profiles}
\label{sec:pumulator}

The reference is the supplied \cmd{pumulator\_infra\_all} source snapshot at
commit \cmd{8fd8bd6f}. The inherited equations are implemented by
\cmd{set\_pum\_latencies()} in
\path{simulator/ramulator2/src/dram/impl/DDR5-PUM.cpp}; the matching defaults
are in \path{simulator/ramulator2/configs/pum/ddr5/simdram.yaml}.
With $t_{\mathrm{RAS}}=\param{nRAS}\,t_{\mathrm{CK}}$ and
$t_{\mathrm{RP}}=\param{nRP}\,t_{\mathrm{CK}}$, they give:
\begin{align*}
 n_{\mathrm{AP}} &= \left\lceil
       \frac{t_{\mathrm{RAS}}+t_{\mathrm{RP}}}{t_{\mathrm{CK}}}
       \right\rceil,\\
 n_{\mathrm{AAP,split}} &= \left\lceil
       \frac{t_{\mathrm{RAS}}+t_{\mathrm{DEC}}+t_{\mathrm{RP}}}
            {t_{\mathrm{CK}}}\right\rceil,\\
 n_{\mathrm{AAP,two\text{-}activate}} &= \left\lceil
       \frac{2t_{\mathrm{RAS}}+t_{\mathrm{RP}}}{t_{\mathrm{CK}}}
       \right\rceil.
\end{align*}
The default maps \param{nMRA} to AP and \param{nCOPY} to split-decoder AAP.
The two-activate profile (Pumulator's \cmd{AAPN}, configuration
\cmd{aap\_variant: naive}) is an alternative sensitivity assumption when a
split decoder is unavailable; it sets
\param{nCOPY}$=2\param{nRAS}+\param{nRP}$ and derives an unoverridden
\param{nNOT} from that copy cost. Neither profile establishes a measured
upper or lower bound on DDRAM-O latency.

Revision M1-P0 used the legacy Pumulator estimate
$\param{nRAS}+\param{nRP}+\lceil\param{nRP}/4\rceil$, also found in
\path{polypum/polypum.py}. Revision \timingrevision{} replaces its
precharge-proportional margin with the graph-based model's explicit 4 ns
decoder delay. This avoids making decoder delay implicitly track the
precharge duration. Provenance here refers to the supplied implementation,
not an independently verified circuit measurement or paper result.

\subsection{Provisional timing ledger and rationale}

Table~\ref{tab:timing-ledger} records why each placeholder exists and what
evidence is required to retire it. The defaults are modeling hypotheses, not
predictions of a fabricated DDRAM-O device.

\begin{table*}[t]
\centering
\caption{Provisional timing ledger, revision \timingrevision{}.}
\label{tab:timing-ledger}
\rowcolors{2}{DDRAMGrey}{white}
\begin{tabularx}{\textwidth}{@{}p{0.12\textwidth}p{0.19\textwidth}X X@{}}
\toprule
Item & Current rule & Reason for the choice & Replacement evidence \\
\midrule
\param{nXACTO} & 8 cycles & Matches the DDRAM-P cross-view placeholder, giving initial A/B comparisons the same assumed coupling delay. It is deliberately independent of the DDR5 speed bin. & Worst-case SPICE separation between opposite-view sensing starts, converted with $\lceil t/t_{CK}\rceil$. \\
\param{nMRA} & $\param{nRAS}+\param{nRP}$ & Inherited Pumulator AP cost: one complete active/restore interval plus local precharge. No extra precharge is charged after this atomic primitive. & Simulated multi-row sensing, restoration, and precharge completion time for each supported row multiplicity. \\
\param{nCOPY} & $\param{nRAS}+\param{nRP}+n_{\mathrm{DEC}}$ & Inherited Pumulator split-decoder AAP cost, with a separately configurable 4 ns decoder delay. Derived from DDR5 base timings, independently of any MRA-specific override. & Confirmation of split-decoder support and an end-to-end source-to-destination copy waveform, including safe local-view reuse. \\
\param{nNOT} & $\param{nCOPY}$ & Inherited Pumulator SIMDRAM assumption: one AAP into a DCC row per inverter. Equality is conditional on that circuit path, not a general inversion latency. & Dedicated inversion-path latency, or evidence that the copy and inversion paths are identical. \\
MRA CA occupancy & 2 cycles & Conservatively matches DDR5 address-bearing activation commands while the final multi-row operand encoding is unspecified. & Final command encoding and command/address bus protocol. \\
MRA activation weight & 1 event & Lets existing \param{nRRD}/\param{nFAW} machinery limit MRA immediately, without claiming that MRA draws ordinary-ACT current. & Calibrated current envelope and either a weighted activation budget or dedicated MRA power constraint. \\
\bottomrule
\end{tabularx}
\end{table*}

The implementation must resolve the DDR5 preset and all base-timing overrides
before deriving these values. Explicit primitive overrides take precedence
and must be positive integer cycles. An unoverridden \param{nNOT} follows
the final resolved \param{nCOPY}; overriding \param{nMRA} alone must not
change \param{nCOPY}.
For the current \cmd{DDR5\_3200AN} preset, \param{nRAS}$=52$ and
\param{nRP}$=24$ with $t_{\mathrm{CK}}=0.625$ ns, yielding
$n_{\mathrm{DEC}}=7$, \param{nMRA}$=76$ cycles (47.5 ns), and
\param{nCOPY}=\param{nNOT}$=83$ cycles (51.875 ns) under the split profile;
the two-activate copy profile instead gives 128 cycles (80 ns).
\param{nXACTO} remains 8 cycles. The old M1-P0 copy value was 82 cycles.
The reserved \param{nCOPY} and \param{nNOT} values do not yet gate an issued
command and therefore must not be interpreted as measured operation latency in
Milestone 1 results.

\subsection{Applicability limits}

Pumulator does not establish DDRAM-O's orthogonal cross-view coupling delay,
intersection behavior, MRA command encoding, or activation-current weight.
Thus \param{nXACTO}, the two-cycle MRA CA occupancy, and the one-event MRA
weight in \param{nRRD}/\param{nFAW} remain DDRAM-specific provisional
assumptions, not Pumulator-inherited timings. A one-event weight does not
limit MRA to four simultaneously activated rows; it approximates its demand
on the shared activation-power budget.

Other Pumulator substrates must not be mixed into the default SIMDRAM profile.
PRADA's 5 ns offset cancellation and 35 ns inversion, ELP2IM's
pseudo-precharge and shortened restoration, and TRDRAM's shift costs require
their respective circuitry. Short-lived intermediate values alone do not
justify reducing the full restore interval of the current atomic MRA.

\subsection{Activation-like command class}

For cross-view constraints, the activation-like class is

\begin{align*}
\mathcal{A}_{A} &= \{\cmd{ACT\_A},\cmd{MRA\_A}\},\\
\mathcal{A}_{B} &= \{\cmd{ACT\_B},\cmd{MRA\_B}\}.
\end{align*}

Future copy and inversion commands must join the class if their physical
implementation begins with activation of their selected view.

\subsection{Required timing relationships}

\begin{table*}[t]
\centering
\caption{Normative timing relationships.}
\label{tab:constraints}
\begin{tabularx}{\textwidth}{@{}p{0.22\textwidth}p{0.22\textwidth}p{0.13\textwidth}X@{}}
\toprule
Preceding & Following & Delay & Scope \\
\midrule
\rowcolor{DDRAMGrey}$\mathcal{A}_{A}$ & $\mathcal{A}_{B}$ & \param{nXACTO} & Same bank/subarray, cross-view \\
$\mathcal{A}_{B}$ & $\mathcal{A}_{A}$ & \param{nXACTO} & Same bank/subarray, cross-view \\
\rowcolor{DDRAMGrey}\cmd{MRA\_A} & Any \viewA{} row command & \param{nMRA} & Local view \\
\cmd{MRA\_B} & Any \viewB{} row command & \param{nMRA} & Local view \\
\rowcolor{DDRAMGrey}\cmd{ACT\_A/B} & Local \cmd{RD}/\cmd{WR} & \param{nRCD} & Inherited DDR5 behavior \\
\cmd{ACT\_A/B} & Local \cmd{PRE} & \param{nRAS} & Inherited DDR5 behavior \\
\rowcolor{DDRAMGrey}\cmd{PRE\_A/B} & Local activation-like & \param{nRP} & Local view \\
\cmd{RD\_A/B}, \cmd{WR\_A/B} & Column commands & DDR5 column constraints & Shared bank/channel data path \\
\bottomrule
\end{tabularx}
\end{table*}

An \cmd{MRA} command does not occupy the data bus. Opposite-view row-path
commands may overlap its local occupancy after satisfying \param{nXACTO}, but
ordinary column commands remain subject to shared DDR5 constraints.

\subsection{Activation power window}

\provisional{} Each \cmd{MRA} counts as one activation event for \param{nFAW}
and the appropriate \param{nRRD} family. This is a bookkeeping approximation,
not a claim that a multi-row activation draws the current of one ordinary
activation. A later calibrated model may introduce a weighted activation
budget or a dedicated constraint.

\section{Prerequisites and State Transitions}

\subsection{Conventional requests}

For a request assigned to view $V$:

\begin{enumerate}
  \item If the target row is open in $V$, issue the terminal column command.
  \item If $V$ is closed, issue \cmd{ACT\_$V$}, then the column command.
  \item If another row is open in $V$, issue \cmd{PRE\_$V$},
        \cmd{ACT\_$V$}, and then the column command.
  \item State in the other view is neither a prerequisite nor a victim.
\end{enumerate}

\subsection{Atomic multi-row activation}

For \cmd{MRA\_$V$}:

\begin{enumerate}
  \item If $V$ is open, issue \cmd{PRE\_$V$} first.
  \item The other view need not be closed.
  \item On issue, set \cmd{busy\_until[$V$]} to the current cycle plus
        \param{nMRA}.
  \item Leave \cmd{active\_row[$V$]} closed.
  \item Do not change the other view's active row or busy deadline.
\end{enumerate}

The controller must preserve an operation's selected view after its first
prerequisite issues, following the same commitment principle used for DDRAM-P.

\section{Layout-Aware Controller Integration}
\label{sec:controller-layout}

The initial controller supports host traffic pinned to A and explicit,
timing-only MRA requests pinned to B. Layout tags and identity/version are
validated, but full operand placement, functional arithmetic, and conversion
commands remain future work. The timing ledger remains \timingrevision{}.

\subsection{Request contract and ownership}

Keep the semantic operation, layout requirement, and current device command
separate. A PuD request descriptor must identify its operation, source and
destination regions, layout identity/version, element width, bit order, lane
mapping, and operand/temporary placement. A layout registry or compiler-owned
descriptor may hold this information; each queued primitive need only carry
a stable reference plus its decoded target and required view. The exact
storage/API for full compound operations remains future work.

The implemented timing-level leaf is \cmd{Request::Type::PudMRA}. It carries
\cmd{required\_view=B}, \cmd{pud\_layout=BitSerial}, nonnegative
\cmd{layout\_id} and \cmd{layout\_version}, and an already-decoded, concrete
B-local address vector. It bypasses the host linear-address mapper so B
coordinates cannot be reinterpreted as A coordinates. Missing/incompatible
tags, invalid coordinates, and an attempted version change while that layout
identity has outstanding PuD leaves are rejected. Identity/version validation
does not establish the physical placement or correctness of operand values.

Successful admission returns a shared \cmd{pud\_completion} token on the
submitted request. An optional \cmd{pud\_dependency} token from an earlier
request in the same controller gates a consumer until service completion.
The initial API supports one explicit predecessor; a compiler/graph executor
will supply more general dependency joins later. Host consumers may also
carry this token. Forwarding/coalescing cannot bypass it; a dependency-blocked
alias of a buffered write applies admission backpressure until a safe retry.
The producer must still describe logical dependencies explicitly; the
timing-only controller does not infer shared-cell conflicts from layout IDs.

The layout/operation mapping produces an immutable \emph{required view} for
each primitive in the initial integration: A for ordinary host reads/writes,
B for bit-serial PuD leaves. This requirement is distinct from the existing
\cmd{selected\_view} and \cmd{view\_committed} execution fields. An explicit
view-local internal command fixes its own view and must agree with any supplied
layout requirement. Unknown or contradictory PuD metadata is an error;
temporary resource contention is a wait condition. Missing PuD metadata must
never silently fall back to the host-access default.

The frontend/PuD compiler supplies the layout contract and operation graph.
The address mapper resolves targets in the required view's coordinate system.
The controller validates the contract and maps commands; the scheduler ranks
eligible requests. The DRAM device remains responsible for prerequisites,
state transitions, and timing checks.

\subsection{Queue organization and arbitration}
\label{sec:controller-queues}

Queues are local to one memory channel; there is no PuD queue per bank or per
view. The controller uses five schedulable request buffers and two
completion-only deques:

\begin{center}
\small
\begin{tabularx}{\columnwidth}{@{}p{0.31\columnwidth}X@{}}
\toprule
Structure & Purpose \\
\midrule
\cmd{read\_buffer} & Admitted host reads. \\
\cmd{write\_buffer} & Admitted host writes; existing forwarding and
coalescing rules still apply. \\
\cmd{pud\_buffer} & Admitted timing-level PuD leaves for all banks in the
channel. Its capacity equals \cmd{read\_buffer\_size}, 32 by default. \\
\cmd{priority\_buffer} & Controller-issued maintenance and explicit direct
commands. MRA is not legal on this path. \\
\cmd{active\_buffer} & Host requests that have issued an opening prerequisite
and must finish their activation-to-access sequence. \\
\cmd{pending} & Issued reads waiting for their data-return deadline; not
schedulable. \\
\cmd{pud\_pending} & Issued MRA leaves waiting for internal service completion;
not schedulable. \\
\bottomrule
\end{tabularx}
\end{center}

The buffers are insertion-ordered lists, but the schedulable demand and PuD
buffers are not strict FIFOs. In each controller cycle, at most one external
DRAM command is selected in the following class order:

\begin{equation}
  \text{active} \;>\; \text{priority} \;>\; \text{PuD}
  \;>\; \text{host read/write}.
\end{equation}

The active class preserves progress for a request whose opening sequence has
already begun. The priority class examines only its head. Consequently, a
nonempty but timing-blocked priority buffer prevents PuD and host arbitration,
although a ready active request may still issue first. Read/write selection
retains the existing watermark-controlled write-drain mode.

Within a demand or PuD buffer, FR-FCFS prepares every request, removes
ineligible entries from consideration, prefers a timing-ready command, and
breaks equal-readiness ties by request arrival time. Eligibility includes an
explicit dependency being complete, expiration of any retry delay, managed
refresh permission, and protection of in-flight active sequences. Therefore,
an older blocked PuD leaf does not prevent an independent younger ready leaf
from issuing. If no PuD leaf is issuable, arbitration proceeds to host
traffic. \cmd{FRFCFS-ViewAware} additionally protects legal view-local row
hits from destructive closes; an atomic MRA has no persistent row-hit state,
so PuD-leaf ordering normally reduces to readiness and age.

Admission, DRAM readiness, issue, and completion are distinct times. Let
$t_a$ be queue admission, $t_r$ the earliest cycle permitted by DRAM timing,
$t_i$ actual command issue, and $t_c$ service completion. They obey

\begin{equation}
  t_i \ge \max(t_a,t_r), \qquad
  t_c = t_i + \param{nMRA}
\end{equation}

for the implemented fixed-latency MRA leaf. Scheduler competition may make the
first inequality strict. A trace timing arrow records the constraint that
determines $t_r$; it does not assert that the target request was already
queued. Visible distance after an arrow's minimum endpoint may therefore be
admission slack or arbitration slack rather than an unmodelled device delay.

\subsection{PuD queue lifecycle}

Admission through \cmd{send\_pud()} validates the B-local address, required
view, bit-serial layout, identity/version, and predecessor token before
appending the request to \cmd{pud\_buffer}. A full buffer applies backpressure.
Every scheduler scan then derives the request's current command from the
device state:

\begin{equation}
\begin{aligned}
  \viewB\text{ closed:} && \cmd{PudMRA} &\rightarrow \cmd{MRA\_B},\\
  \viewB\text{ open:}   && \cmd{PudMRA} &\rightarrow \cmd{PRE\_B}
      \rightarrow \cmd{MRA\_B}.
\end{aligned}
\end{equation}

An issued \cmd{PRE\_B} is a prerequisite of the same request. The request
stays in \cmd{pud\_buffer}, retains its decoded coordinates and committed
view, and is reconsidered until \cmd{MRA\_B} becomes ready after \param{nRP}.
It does not enter \cmd{active\_buffer}: PRE is closing, while MRA is an atomic
operation that intentionally leaves no conventional row open.

When the final \cmd{MRA\_B} issues at $t_i$, the request leaves the schedulable
buffer and enters \cmd{pud\_pending} with departure time
$t_i+\param{nMRA}$. Controller view-selection ownership ends at MRA issue, but
the device timing ledger prevents same-view use during the remaining service
interval. The other view may progress once \param{nXACTO} and all shared
constraints permit it. At departure, the controller marks the shared
completion token complete, releases the outstanding layout reference, updates
the served counter, and invokes the callback. A consumer carrying that token
is ineligible until this point, not merely until MRA issue. Fixed \param{nMRA}
makes \cmd{pud\_pending} completion-ordered; this deque is never examined by
the command scheduler.

PuD counters \cmd{num\_pud\_reqs} and \cmd{num\_pud\_reqs\_served} count
admitted and service-completed leaves, respectively. They do not contribute
to read latency or maintenance counts. Direct priority-queue MRA submission
is rejected because it would bypass metadata validation, dependency tracking,
and delayed service completion. A-view PuD remains deferred.

\paragraph{Cross-view trace example.}
The development trace intentionally separates phases. Its first
\cmd{PRE\_A} is admitted only after earlier reads drain, although
\param{nRAS} made it timing-ready sooner. Likewise, the second \cmd{MRA\_B}
is admitted after an explicit \param{nRP} delay, long after the preceding
MRA's \param{nMRA} bound has expired. Those gaps are admission slack. The
subsequent \cmd{MRA\_B}$\rightarrow$\cmd{ACT\_A} gap is instead tight device
timing: \cmd{ACT\_A} issues exactly \param{nXACTO} cycles after MRA. The
example must not be interpreted as evidence that the scheduler held an
already-ready PRE or MRA for the entire visible gap.

\subsection{Command mapping and scheduling pipeline}

The initial mappings are:
\begin{itemize}
  \item Host read/write: \cmd{RD\_A}/\cmd{WR\_A}, with local
        \cmd{PRE\_A}/\cmd{ACT\_A} prerequisites as needed.
  \item Bit-serial multi-row primitive: \cmd{MRA\_B}, with
        \cmd{PRE\_B} first if its local view is open.
  \item Explicit B-local read/write: \cmd{RD\_B}/\cmd{WR\_B} using B-local
        coordinates and a compatible layout descriptor.
  \item Future COPY/NOT leaves in the initial PuD integration: use
        \cmd{COPY\_B}/\cmd{NOT\_B} when those commands and operand semantics
        are implemented. A-view PuD leaves remain deferred. The
        reserved timing parameters alone do not implement these operations.
\end{itemize}

Extend the existing controller-level \cmd{prepare\_request()} path: validate
the descriptor, resolve the legal target/view, set \cmd{final\_command}, and
then ask the device for the next prerequisite. Preparation during scheduler
scans must not issue commands, allocate compute rows, mutate layout ownership,
or reserve a view. Validate and decode immutable metadata at admission where
possible; readiness may change on subsequent scans.

Reuse the \cmd{IViewSelection} boundary for mapping and commitment, with a
DDRAM-O layout-constrained policy. Do not simply enable DDRAM-P's
\cmd{TwoView}/\cmd{QueueAware} policies for DDRAM-O: those policies assume
either view can serve a data request and treat equal row IDs as the same
physical row. DDRAM-O uses different coordinate namespaces and layout-pinned
targets. No opposite-view row hit, free view, queue preference, or activation
NACK may redirect a pinned request to another view.

FR-FCFS continues to rank prepared, dependency-ready requests by readiness
and arrival time. A row-hit-preferring scheduler may favor hits only in a
request's legal view. It must not choose the layout or translate coordinates.
An atomic MRA does not have a persistent open-row hit; its eligibility follows
its local prerequisites and timing. A blocked B operation may coexist with
independent A progress, subject to shared timing and logical dependencies.

\subsection{Commitment, dependencies, and completion}

Keep the required view and decoded target stable throughout execution.
After the first successfully issued prerequisite, including PRE, commit the
selected view and retain it across queue moves and retries. An activation
NACK does not commit a previously unstarted request; if PRE already issued,
the earlier commitment remains. Any retry uses the same legal view and
coordinates. Resource ownership must preserve progress without allowing an
incompatible request to invalidate the in-flight prerequisite sequence.

For the first integration, freeze layout/placement while dependent requests
are outstanding. Future relocation or conversion must drain or order affected
requests and update the layout version explicitly. The operation graph must
order conflicting host/PuD accesses and operand producers/consumers; independent
view state is not evidence that shared-cell data dependencies are absent.
Milestone 1 still omits electrical intersection hazards and functional cell
mapping, so timing-only overlap tests cannot establish functional correctness.

Issuing an MRA leaf starts its service interval. A dependent leaf or a parent
PuD completion notification must wait for that service to finish, at least
\param{nMRA} cycles after issue, in addition to graph dependencies. Ordinary
read completion retains the existing read-latency path. Do not apply a read
data-bus return latency to MRA or report a complete compound operation when
its terminal primitive merely issues. COPY/NOT completion will use their
own service intervals once implemented.

\subsection{Incremental integration and validation}

Start with host-only A mapping on DDRAM-O and compare command histories and
latency with DDR5. Next inject explicit, layout-pinned B-local MRA requests
to validate mapping, PRE/MRA commitment, service completion, and A/B overlap.
These are timing-only primitives, not functional arithmetic. Compound
bit-serial operations follow once operand placement and the functional model
are available. Required first controller tests are:
\begin{enumerate}
  \item Host read/write traffic uses A and matches DDR5 under equivalent
        presets, maintenance, row policy, and workload.
  \item A B-required PuD request remains on B even when A is idle or has an
        equal numeric row ID open; a host request likewise stays on A.
  \item B-local target coordinates survive prerequisite resolution and queue
        movement; an open B view causes PRE\_B before MRA\_B without closing A.
  \item Repeated preparation, an issued PRE, and activation retries cannot
        change the required view, descriptor version, or decoded target.
  \item Missing, incompatible, or contradictory PuD layout metadata is rejected;
        occupied resources wait without triggering an implicit conversion.
        A-view PuD requests are rejected as outside the initial controller
        scope, while ordinary A-view host accesses remain supported.
  \item Dependent operations and completion wait for MRA's full service
        interval; independent A work can overlap when shared constraints allow.
  \item Conventional DDR5 and DDRAM-P controller regressions remain unchanged.
\end{enumerate}

\section{Configuration Interface}

The Python DRAM specification must accept explicit overrides. A representative
configuration is:

\begin{small}
\begin{verbatim}
ramulator.dram.DDRAM_O(
    org_preset="DDR5_16Gb_x8",
    timing_preset="DDR5_3200AN",
    rank=1,
    nXACTO=8,
    nMRA=76,
    nCOPY=83,
    nNOT=83,
)
\end{verbatim}
\end{small}

The example explicitly supplies the \timingrevision{} split-profile values;
omitting those primitive overrides now derives the same defaults.
Derived defaults must be resolved after
the inherited DDR5 preset and base overrides so that changing \param{nRAS}
or \param{nRP} changes the placeholder consistently. Decoder delay must
remain a configurable physical-time assumption, initially 4 ns. The implemented
constructor accepts \cmd{tDEC\_ns=4.0} and
\cmd{copy\_timing\_profile="pumulator-split"}; the alternative is
\cmd{"pumulator-two-activate"}. Both use ceiling rounding. Configuration
metadata in \cmd{ddram\_timing\_model} records the ledger revision, copy
profile, decoder delay, presets, explicit overrides, and intersection flag.
Every experiment result must record:

\begin{itemize}
  \item the resolved values of \param{nXACTO}, \param{nMRA},
        \param{nCOPY}, and \param{nNOT};
  \item the DDR5 organization and timing presets;
  \item the copy profile (split or two-activate), decoder delay in ns,
        clock period, rounding rule (ceiling), and any explicit overrides;
  \item a timing-model revision string, initially
        \cmd{\timingrevision}; and
  \item whether intersection modeling is enabled, initially \cmd{false}.
\end{itemize}

Invalid negative or zero DDRAM-O-specific timing values must be rejected after
default resolution. The decoder-delay input and clock period must be finite
and positive.

\paragraph{Implementation alignment.}
The simulator implements \timingrevision{} defaults, both copy profiles, and
PuD-default derivation after explicit base/primitive overrides, locally to
DDRAM-O. An unoverridden NOT follows the final COPY value; MRA overrides do
not change COPY. Conventional DDR5 parameter resolution is unchanged.
COPY/NOT remain reserved costs rather than issued commands. Runs must record
actual resolved values and distinguish timing-only leaves from functional
compound operations.

\section{Trace and Statistics Requirements}

The existing trace recorder must represent every DDRAM-O command without
changing conventional trace semantics. Each event should expose:

\begin{itemize}
  \item command name and issue cycle;
  \item selected view;
  \item bank/subarray and representative row address;
  \item whether it is a conventional or PuD command;
  \item the timing constraint that determined readiness; and
  \item the resolved timing value used by that constraint.
\end{itemize}

The initial statistics should include counts of view-A and view-B activations,
multi-row activations by view, cycles each view is busy, cross-view overlap
cycles, and stalls attributed to \param{nXACTO} or \param{nMRA}.

\section{Verification Plan}

Milestone 1 is complete only when the following tests pass.

\begin{table*}[t]
\centering
\caption{Acceptance tests.}
\label{tab:tests}
\begin{tabularx}{\textwidth}{@{}p{0.06\textwidth}p{0.36\textwidth}X@{}}
\toprule
ID & Scenario & Required observation \\
\midrule
\rowcolor{DDRAMGrey}S1 & Open one row through each view. & Both rows remain open independently. \\
S2 & Precharge one view. & Only that view closes. \\
\rowcolor{DDRAMGrey}S3 & Open equal view-local line indices in A and B. & Both activations are accepted; equality of namespace-local labels creates no Milestone 1 intersection rule. \\
T1 & Issue opposite-view activations. & Start times differ by at least \param{nXACTO}. \\
\rowcolor{DDRAMGrey}T2 & Repeat T1 with an override. & The observed separation follows the configured value exactly when no stricter constraint applies. \\
T3 & Issue \cmd{MRA\_A}, then an A operation. & The A operation waits at least \param{nMRA}. \\
\rowcolor{DDRAMGrey}T4 & Issue \cmd{MRA\_A}, then a B operation. & B may start after cross-view constraints without waiting for the full \param{nMRA}. \\
T5 & Exercise the four-activation window. & MRA contributes one provisional activation event. \\
\rowcolor{DDRAMGrey}P1 & Issue MRA while its local view is open. & A local precharge is issued first; the other view remains open. \\
P2 & Issue rank or same-bank maintenance. & Both views close before inherited maintenance executes. \\
\rowcolor{DDRAMGrey}C1 & Run conventional DDR5 tests. & Command translation, timing, and state behavior are unchanged. \\
C2 & Record a mixed A/B/MRA sequence. & Trace contains view identity and the expected timing dependencies. \\
\rowcolor{DDRAMGrey}D1 & Resolve the split profile at DDR5\_3200AN. & MRA is 76 cycles, COPY and unoverridden NOT are 83 cycles; the decoder delay is rounded up from 4 ns to 7 cycles. \\
D2 & Change speed bin, decoder delay, or base timings. & Defaults follow the physical-time formulas after base overrides, with ceiling conversion; two-activate COPY uses $2\param{nRAS}+\param{nRP}$. \\
\rowcolor{DDRAMGrey}D3 & Override MRA, COPY, and NOT independently. & MRA alone does not alter COPY; an unoverridden NOT follows resolved COPY; explicit NOT wins. Invalid values are rejected. \\
\bottomrule
\end{tabularx}
\end{table*}

Every timing test should independently replay the issued history through the
device prerequisite and timing checks, as the DDRAM-P suite already does.

\section{Calibration and Evolution}

The command interface is intended to remain stable while the physical model
becomes more detailed.

\subsection{SPICE calibration}

The first calibration step replaces \param{nXACTO} using the minimum safe
separation between opposite-view sensing events under the selected voltage,
temperature, process corner, and row multiplicity. The report should provide
the time in nanoseconds and the conversion rule

\begin{equation}
  \param{nXACTO} = \left\lceil
    \frac{t_{\mathrm{XACTO}}}{t_{\mathrm{CK}}}
  \right\rceil.
\end{equation}

Multi-row activation simulations should subsequently calibrate local-view
occupancy, restoration time, current limits, and whether MRA should carry more
than one activation unit in the power window.

Each calibration update must increment the timing-ledger revision, retain the
superseded defaults in version history, and record the circuit configuration,
process/voltage/temperature corner, row multiplicity, time in nanoseconds,
clock period, rounding rule, and resulting cycle count. This makes results
produced with provisional and calibrated models distinguishable.

\subsection{Intersection-aware extension}

A later revision may add intersection identity, masking, and simultaneous
sensing restrictions. It should extend readiness checks and functional data
mapping, not replace view-specific command names or provisional timing
parameters. Experiments must report which revision they use.

\subsection{Functional PuD extension}

Milestone 2 will attach explicit operand and destination rows to compound PuD
requests. A functional shadow array will map the orthogonal coordinate systems
onto shared cells and validate majority, addition, and multiplication. The
atomic \cmd{MRA\_B} command defined here is the initial multi-row timing
leaf of those operation graphs; A-view PuD remains a deferred extension.

\subsection{Intermediate-result cost accounting}

Pumulator's \cmd{Model::build\_simdram()} provides a reproducible baseline
for future graph costs. Its source is:
{\raggedright\noindent
\path{simulator/ramulator2/src/frontend/impl/processor/simpleO3/pum_model.cpp}\par}
Its default MAJ3
sequence copies three operands into compute rows, executes one AP, and
copies one result out. Mapping AP to MRA gives the serialized service cost
\begin{equation}
 n_{\mathrm{MAJ3}}=3\param{nCOPY}+\param{nMRA}+\param{nCOPY}.
\end{equation}
At DDR5\_3200AN under the split profile this is 408 cycles (255~ns),
before extra inversion, layout conversion, movement between views, queueing,
maintenance, and shared-resource contention. This is a primitive-sequence
estimate, not a simulated end-to-end latency or the cost of one MRA alone.
The final precharge is already included in each primitive's occupancy;
an initially open local view still needs its ordinary prerequisite precharge.

Milestone 2 must retain dependency edges and operand lifetimes, track which
values are overwritten by destructive multi-row sensing, and account for
copies that preserve live intermediate values. Copy elision requires proven
residency and liveness; a value in A is not automatically available in B at
zero cost. Fused or free inversion requires a supported physical path and
placement. Pumulator's distance-based spill heuristic can inform experiments
but is not a substitute for an allocator with explicit row capacity.
Schedule the resulting primitives through actual view, command-bus, and
activation constraints; do not divide serialized costs by assumed bank or
subarray parallelism factors to claim DDRAM-O latency.

\section{Implementation Boundary}

The recommended separation of responsibilities is:

\begin{itemize}
  \item \textbf{DRAM specification:} commands, state transitions, prerequisites,
        timing constraints, and parameter resolution;
  \item \textbf{controller:} request preparation, selected-view commitment,
        timing-aware issue, and retirement;
  \item \textbf{PuD compiler (later):} operation decomposition, dependencies,
        operand placement, and temporary rows;
  \item \textbf{scheduler (later):} selection among ready primitive commands;
  \item \textbf{functional model (later):} data-value correctness; and
  \item \textbf{energy model (later):} calibrated per-primitive energy.
\end{itemize}

This boundary preserves conventional DDR5 behavior and prevents placeholder
electrical assumptions from becoming embedded in scheduling policy.

\end{document}
